\documentclass[12pt,a4paper]{article}
\usepackage[dvipsnames]{xcolor}
\usepackage[utf8]{inputenc}
\usepackage{amsmath}
\usepackage{mathtools}
\usepackage{marvosym}
\usepackage{wrapfig}
\usepackage{hyperref}
\usepackage{float}
\usepackage{multicol}
\hypersetup{colorlinks,citecolor=black,filecolor=black,linkcolor=black,urlcolor=black}
\usepackage{pdfpages}
\usepackage{amsfonts}
\usepackage{amssymb}
\usepackage{fancyhdr}
\usepackage{chemfig}
\usepackage{graphicx}
\usepackage[T1]{fontenc}
\usepackage[magyar]{babel}
\usepackage{bm}
\usepackage{tikz, tcolorbox}
\usetikzlibrary{positioning}
\usepackage{verbatim}
\usepackage{amsmath,esint}
\usepackage{setspace}
\usepackage{qtree}
\usepackage{multido}
\usepackage{mathptmx} % Times New Roman betűtípus
\usepackage{listings}
\usepackage{caption}
\usepackage{xparse}
\captionsetup[lstlisting]{labelformat=empty}

\lstset{
    language=Matlab,
    inputencoding=utf8,
    extendedchars=true,
    basicstyle=\ttfamily\small,
    frame=single,
    breaklines=true,
    numbers=left,
    numberstyle=\tiny,
    literate=
        {á}{{\'a}}1
        {é}{{\'e}}1
        {í}{{\'i}}1
        {ó}{{\'o}}1
        {ö}{{\"o}}1
        {ő}{{\H{o}}}1
        {ú}{{\'u}}1
        {ü}{{\"u}}1
        {ű}{{\H{u}}}1
        {Á}{{\'A}}1
        {É}{{\'E}}1
        {Í}{{\'I}}1
        {Ó}{{\'O}}1
        {Ö}{{\"O}}1
        {Ő}{{\H{O}}}1
        {Ú}{{\'U}}1
        {Ü}{{\"U}}1
        {Ű}{{\H{U}}}1
}
\renewcommand{\figureautorefname}{ábra}
\newcommand{\abrara}[1]{\hyperref[#1]{\ref*{#1}.~ábra}}
\NewDocumentCommand{\kep}{O{300px} O{0} m m m}{%
\begin{figure}[H]
    \centering
    \includegraphics[width=#1, angle=#2]{#3}
    \caption{#4}
    \label{fig:#5}
\end{figure}%
}
\newcommand{\Pointilles}[1]{%
\par\nobreak
\noindent\rule{0pt}{1.5\baselineskip}% Provides a larger gap between the
\multido{}{#1}{\noindent\makebox[\linewidth]{\dotfill}\endgraf}
\bigskip% Gap between dots and next paragraph
}
\newcommand{\feladat}[1]{%
\phantomsection
\subsection*{#1. feladat}%
\addcontentsline{toc}{subsection}{#1. feladat}%
}
\usepackage{pgfplots}
\pgfplotsset{height = 10cm, width=15cm,compat=1.9}
\usepackage[left=2cm,right=2cm,top=2cm,bottom=2cm]{geometry}
\setlength{\parindent}{0pt}
\setlength{\parskip}{0em}
\pagestyle{fancy}
\fancyhf{}
\cfoot{\thepage. oldal}

% Másfeles sorköz beállítása
\onehalfspacing

\begin{document}

\begin{titlepage}
\centering

\includegraphics[width=.8\textwidth]{../../Kepek/bme_logo_nagy.jpg}

\vspace{1em}

{
\large
BUDAPESTI MŰSZAKI ÉS GAZDASÁGTUDOMÁNYI EGYETEM\\ VILLAMOSMÉRNÖKI ÉS INFORMATIKAI KAR\\ MESTERSÉGES INTELLIGENCIA ÉS RENDSZERTERVEZÉS TANSZÉK
}

\vspace{6em}

{
\Huge
\textbf{Elosztott adatgyűjtő rendszerek}
}

\vspace{3em}

{
\huge
Intelligens beágyazott rendszerek laboratórium\\

\vspace{.2em}

\large
(BMEVIMIMA21)
}

\vspace{9em}

{
\Large
\textit{Készítette:}\\
\vspace{.5em}
\textbf{Pál Patrik, Gavallér Dávid}\\
\vspace{.5em}
\textbf{AIQTQW, ES02RO}
}

\vspace{12em}

{
\large
BUDAPEST, \today
}
\end{titlepage}

\tableofcontents

\newpage

\section{Bevezetés}

\subsection*{A mérés célja}

A mérés célja az elosztott rendszerű hálózatok és azok problémáinak megismerése gyakorlati példákon keresztül.\\

Az első fő feladatrészben két mitmótot felhasználva vizsgáljuk a szenzor és a gateway által rögzített jelek szinkronizációs problémáit. A feladat során a két mérőegység által rögzített jelek időbeli összehangolását végezzük el időbélyeg-transzformáció segítségével, majd a gateway mintavételi időpontjaiban interpolációval becsüljük a szenzorértékeket.\\

\subsection*{A méréshez használt eszközök}
\begin{itemize}
	\item Mitmótok
	\item Analóg interface kártya
	\item Oszcilloszkóp
	\item Számítógép
	\item Függvénygenerátor
\end{itemize}

\newpage

\section{Mérési feladatok}

\feladat{1}

Az első feladat célja a különálló mérőegységek által rögzített adatok időbeli
összehangolása. A feladat során először beolvassuk és ábrázoljuk a szenzor, az
adatgyűjtő, valamint a szinkronizációs mérések nyers adatait. A szinkronizációs
pontokra illesztett lineáris időbélyeg-transzformáció segítségével az
adatgyűjtő időskáláját a szenzor időskálájára alakítjuk, majd megvizsgáljuk az
így kapott jelek egyezését. Az összehangolt adatokból lineáris interpolációval
becsüljük a szenzorértékeket az adatgyűjtő mintavételi időpontjaiban.
Az elemzést 50~Hz-es, 500~Hz-es, valamint változó amplitúdójú mérésre is
elvégeztük.

\feladat{1.1}

Az első feladat első pontja a mérési összeállítás elkészítése volt. Külső jelként egy függvénygenerátort használtunk, amit a Low Pass Filter bemenetére kötöttünk. Az LPS kimenetét az ellenállásokon keresztül a szenzor és a gateway bemenetére kapcsoltunk, ide csatlakoztatva egy oszcilloszkópot is, a jel valós idejű ellenőrzésének céljából. A gateway-t végül a számítógéphez csatlakozattuk, hogy a méréseket tudjuk rögzíteni és kiértékelni. A DIP kapcsolókat a mérési útmutatóban leírt módon beállítottuk, majd a jelet megvizsgáltuk az oszcilloszkópon és matlabban is. Az összeállítás az alábbi blank. ábrán látható:

%\kep{../../Kepek/20260914_142213.jpg}{A mérési elrendezés.}{meresi-elrendezes}

\feladat{1.2}
A következő feladat a mérési regisztrátum készítése volt. A mérés rögzítéséhez a mérési jegyzőkönyvben leírt programot használtuk, majd a \texttt{preprocessSampleFiles.m} script futtáatása után kapott, egyszerűen feldolgozható fájlokat vizsgáltuk. A mérést elvégeztük 50~Hz-es és 500~Hz-es szinuszos jelekre is, úgy, hogy mérés közben ki és bekapcsoltuk a jelgenerátort. A jelek amplitúdóját mind a kettő esetben 2 V-ra állítottuk, ezzel kitöltve a mérési tartományt. A mérési eredmények 50~Hz-es és 500~Hz-es jelekre az alábbi blank. ábrákon láthatóak:

%y tengely felirat? mit mertunk

\kep{../../Kepek/Nyers_Szezor_50hz.jpg}{Nyers szenzoradatok (50 Hz).}{nyers-szenzoradatok-50hz}

\kep{../../Kepek/Nyers_Szenzor_500hz.jpg}{Nyers szenzoradatok (500 Hz).}{nyers-szenzoradatok-500hz}

\feladat{1.3}
A harmadik feladatban az volt a célunk, hogy a szenzor és a gateway által rögzített jeleket időben összehangoljuk. Az időbeli eltérést az eszközöknél kezdeti időeltérés a két belső óra skáláján és a két órajel mintavételi frekvenciája közötti elrétés is okozhatta. Ennek a megoldására egy olyan lineáris transzformációt kerestünk, ami a gateway időskáláját a szenzor időskálájára alakítja.\\

Ehhez az \texttt{stmpSync} fájlt használtuk fel, amiben szinkronizációs időbélyeg párok szerepelek. Ez azt jelenti, hogy ezek az adatok két eszköz által rögzített, azonos eseményekhez tartozó időbélyegek. A Matlab \texttt{polyfit} függvényével ezekre az időbélyegekre lineáris függvényt illesztettünk, ami az alábbi ábrákon látható 50~Hz-es és 500~Hz-es jelekre:

\kep{../../Kepek/Szinkronizacios_idobelyeg_50Hz.jpg}{Szinkronizációs időbélyegek (50 Hz).}{Szinkronizacios-idobelyeg-50Hz}

\kep{../../Kepek/Szinkronizacios_idobelyeg_500Hz.jpg}{Szinkronizációs időbélyegek (500 Hz).}{Szinkronizacios-idobelyeg-500Hz}

A \texttt{polyfit} függvény az illesztett elsőfokú polinom együtthatóit adja
vissza. A méréshez tartozó szinkronizációs pontokra az illesztés eredménye közelítőleg
\begin{equation}
    \texttt{idobelyegPoli}=[1,\,-14{,}9512]
\end{equation}

Mivel elsőfokú polinomot illesztettünk, az első érték a meredekség, a második pedig a tengelymetszet. Ezért a gateway időbélyegét a következő összefüggéssel alakítottuk át a szenzor időskálájára:

\begin{equation}
    t_{\mathrm{szenzor}} = 1 \cdot t_{\mathrm{gateway}} - 14{,}9512
\end{equation}

A tengelymetszet pedig azt a kezdeti időeltérést írja le, amely a két eszköz időbélyegei között fennáll.

Az illesztési együtthatókat a \texttt{polyval} függvénnyel alkalmaztuk \texttt{gwySampX} minden elemére, így kiszámítva az illesztett polinom által becsült, a szenzor időskálájára átszámított időbélyegeket, amik az alábbi ábrákon láthatóak 50~Hz-es és 500~Hz-es jelek esetén:

\kep{../../Kepek/Idobelyeg_50hz.jpg}{Időbélyeggel összehangolt jelek (50 Hz).}{idobelyeg-50Hz}

\kep{../../Kepek/Idobelyeg_500hz.jpg}{Időbélyeggel összehangolt jelek (500 Hz).}{idobelyeg-500Hz}

Ezt követően egy mérési regisztrátumot készítettünk, amelyben a jel amplitúdóját az indulástól kezdve folyamatosan növeltük, majd elvégeztük rá a korábbi feladatokban leírt összehangolását. Ez alatt a mérés alatt 50~Hz-es frekvencián mértünk. Az alábbi ábrákon látható a nyers mérési regisztrátum, a szinkronizációs időbélyegek, valamint az összehangolt jelek:

\kep{../../Kepek/Nyers_szenzor_valtozo.jpg}{Nyers szenzor jelek változó amplitúdóval.}{nyers-szenzor-valtozo}

\kep{../../Kepek/Szinkronizacios_idobelyeg_valtozo.jpg}{Szinkronizációs időbélyegek változó amplitúdóval.}{szinkronizacios-idobelyeg-valtozo}

\kep{../../Kepek/Idobelyeg_valtozo.jpg}{Időbélyeggel összehangolt jelek változó amplitúdóval.}{idobelyeg-valtozo}

\feladat{1.4}

A negyedik feladatban az előző feladatban bemutatott időbélyeg-transzformációt végeztük el újra, LS-becslés segítségével. Felépítettük az \textbf{X} mátrixot, amelynek első oszlopa a gateway szinkronizációs időbélyegeit, második oszlopa pedig egyeseket tartalmazott. Az \textbf{y} vektor a szenzor szinkronizációs időbélyegeit tartalmazta. A számításhoz a Moore--Penrose-pszeudoinverzt használtuk, amelynek LS-megoldása:

\begin{equation}
    p = (\textbf{X}^T\textbf{X})^{-1}\textbf{X}^T\textbf{y}
\end{equation}

A MATLAB beépített függvényei között megtalálható a \texttt{pinv} függvény, amellyel ez a megoldás SVD-alapú módon számítható. A pszeudoinverzzel összehangolt jelek az alábbi ábrán láthatóak:

\kep{../../Kepek/Pszeudoinverz.jpg}{Pszeudoinverz segítségével összehangolt jelek változó amplitúdóval.}{Pszeudoinverz}

Ezt követően a gateway szinkronizációs időadatait és a gateway mérési időtengelyét $10^6$ másodperccel eltoltuk, majd így ismét elvégeztük a feladatot, és vizsgáltuk a mátrix kondicionáltságát. Az első esetben a kondíciószám $1{,}0367 \cdot 10^7$ volt, az időeltolást követően pedig ez az érték $2{,}6496 \cdot 10^{11}$-re növekedett. Az ábrán ennek ellenére nem látható lényeges eltérés, mivel a lineáris transzformáció konstans tagja kompenzálja az időkoordináta közös eltolását. A nagyobb kondíciószámból azonban következik, hogy a becslés érzékenyebbé válik a mérési és kerekítési hibákra.

A probléma megoldására az időadatokat az illesztés előtt központosítottuk. Ehhez a gateway szinkronizációs időadataiból kivontuk azok átlagát, és ugyanezt az eltolást alkalmaztuk a gateway mérési időpontjaira is. Így a transzformáció eredménye nem változik, miközben a mátrix kondicionáltsága javul.

\feladat{1.5}

Az ötödik feladatban a szenzor mérési eredményeit a gateway mintavételi időpontjaiban kellett meghatározni lineáris interpoláció segítségével. Először a szinkronizációs pontokból lineáris időbélyeg-transzformációt hajtottunk végre, a korábbi feladatokban leírt módszer segítségével. Ezután a gateway minden mintavételi időpontját átszámítottuk a szenzor időrendszerébe. Az \texttt{interp1} függvény ezekben az időpontokban a szenzor két szomszédos mintája közötti egyenest használva becsülte meg a szenzor mérési értékét, az alábbi összefüggés szerint:

\[
\hat{s}(t) =
s(t_i) + \frac{t-t_i}{t_{i+1}-t_i}
\bigl(s(t_{i+1})-s(t_i)\bigr),
\qquad t_i \leq t \leq t_{i+1}.
\]

A kapott interpolált szenzorértékeket a gateway ugyanazon időpontokban mért értékeivel hasonlítottuk össze. Az összehasonlításhoz az átlagos négyzetes eltérést is kiszámítottuk:

\[
\mathrm{MSE} =
\frac{1}{N}\sum_{i=1}^{N}
\left(y_{\mathrm{gateway},i}
-\hat{y}_{\mathrm{szenzor},i}\right)^2.
\]

A programot 50~Hz-es és 500~Hz-es mérésre is lefuttattuk. Az 50~Hz-es jel esetén az MSE értéke $0,7255$, 500~Hz-es jel esetén pedig $1,9439$ volt. Az 500~Hz nagyobb változási sebessége volt a nagyobb hiba oka a második esetben. A mérési eredmények az alábbi ábrákon láthatóak, kellő nagyítással a jobb vizualizáció érdekében:

\kep{../../Kepek/Linear_interp_50hz.jpg}{Lineáris interpoláció 50~Hz-en.}{linear-interp-50hz}

\kep{../../Kepek/Linear_interp_500hz.jpg}{Lineáris interpoláció 500~Hz-en.}{linear-interp-500hz}

\newpage
\section*{Függelék}
\subsection*{MATLAB kódok}

\lstinputlisting[
   caption={Az 1-4 mérés MATLAB-kódja},
   label={lst:meres1-matlab}
]{../../MATLAB/meres1_1234.m}


\lstinputlisting[
   caption={Az 5. mérés MATLAB-kódja},
   label={lst:meres5-matlab}
]{../../MATLAB/meres1_5.m}


\end{document}
